\documentclass[10pt,a4paper]{article}
\usepackage[utf8]{inputenc}
\usepackage[a4paper,left=3cm,right=2cm]{geometry}
\usepackage{amsmath}
\usepackage[thmmarks, amsmath, thref]{ntheorem}
\usepackage{amsfonts}
\usepackage{amssymb}
\usepackage{fancyhdr}
\pagestyle{fancy}
\usepackage{anysize}
\usepackage{graphicx}
\usepackage{float}
\usepackage{hyperref}
\usepackage{enumitem}
\newtheorem{theorem}{Theorem}
\theoremsymbol{\ensuremath{\blacksquare}}
\newtheorem*{solution}{Solución:}

\usepackage{xcolor}

% Margenes y formalidades

\textwidth 6.25in % Margen Derecho y ancho del texto
\textheight 8.5in
\oddsidemargin 0in % Margen Izquierdo
\headheight 0in % Largo del texto
\leftmargin 1in
\parindent 0em % Salto entre la indentación (Sangría)https://www.overleaf.com/4519127669sdpptkhvsbfn
\topmargin -0.5in % Margen Superior
\parskip 0.5ex % Salto entre parrafos
\baselineskip 1pt



\lhead{
\includegraphics[scale=0.11]{figuras/logoUSM.png}} % texto izquierda de la cabecera
%\chead{}                                      % texto centro de la cabecera
\rhead{Universidad Técnica Federico Santa María\\Ingeniería Comercial } % número de página a la derecha
\renewcommand{\headrulewidth}{0.4pt} % grosor de la línea de la cabecera
\renewcommand{\footrulewidth}{0.4pt} % grosor de la línea del pie

\begin{document}
\vspace{5cm}



\begin{center}
\begin{large}
\textbf{\\Guía 5: Econometría ICS-294}
\end{large}\\
\vspace{0.2cm}
I Semestre 2025\\
Profesora: Jocelyn Tapia Stefanoni\\
\end{center}


Un analista político de la oposición está estudiando si los gastos
de campaña efectivamente afectan los resultados de las elecciones.
Dado que desde el gobierno se trabaja en un proyecto de ley que establece
que todos los gastos de campaña deben ser financiados con fondos públicos,
los resultados del estudio pueden ser muy relevantes para la aprobación
de este proyecto.

El modelo general que el analista está estimando es el siguiente:

\begin{equation}
voteA=\beta_{0}+\beta_{1}lexpendA+\beta_{2}lexpendB+\beta_{3}prtystrA+\mu\label{eq:Main equation}
\end{equation}

donde $voteA$ es el porcentaje de votos recibidos por el candidato
A, $expendA$ y $expendB$ son los gastos del campaña del candidato
A y del candidato B, respectivamente y $prtystrA$ es una medida de
la popularidad del partido del candidato A, medida como el porcentaje
de votos que obtuvo el partido del candidato A en la elección presidencial
más reciente. En la ecuación \eqref{eq:Main equation}, $lexpendA$
corresponde al logaritmo natural de los gastos de campaña del candidato
A . %Estadística descriptiva para las variables del estudio se presentan en el cuadro \ref{tab:one}. 

% \begin{table}[!htbp] \centering    \caption{\label{tab:one}Estad\'istica Descriptiva}    \label{}  \begin{tabular}{@{\extracolsep{5pt}}lccccccc}  \\[-1.8ex]\hline  \hline \\[-1.8ex]  Statistic & \multicolumn{1}{c}{N} & \multicolumn{1}{c}{Mean} & \multicolumn{1}{c}{St. Dev.} & \multicolumn{1}{c}{Min} & \multicolumn{1}{c}{Pctl(25)} & \multicolumn{1}{c}{Pctl(75)} & \multicolumn{1}{c}{Max} \\  \hline \\[-1.8ex]  voteA & 173 & 50.503 & 16.785 & 16 & 36 & 65 & 84 \\  expendA & 173 & 310.611 & 280.986 & 0.300 & 81.630 & 457.410 & 1,470.670 \\  expendB & 173 & 305.088 & 306.278 & 0.930 & 60.050 & 450.720 & 1,548.190 \\  prtystrA & 173 & 49.757 & 9.984 & 22 & 44 & 56 & 71 \\  lexpendA & 173 & 5.026 & 1.602 & $-$1.197 & 4.402 & 6.126 & 7.293 \\  lexpendB & 173 & 4.944 & 1.571 & $-$0.073 & 4.095 & 6.111 & 7.345 \\  \hline \\[-1.8ex]  \end{tabular}  \end{table} 


Los resultados de las estimaciones preliminares que ha llevado a cabo
el analista son los siguientes:

\begin{table}[ht] \centering    \caption{\label{tab:two}Estimación por MCO} \label{Resultados MCO} \begin{tabular}{lccccc} \hline \hline                                         & \multicolumn{5}{c}{Dependent variable: vote A}                                     \\ \hline                                         & 1               & 2              & 3             & 4              & 5              \\ \cline{2-6}  lexpendA                                & 6.506***        &                &               & 6.342***       & 6.083***       \\                                         & (0.628)         &                &               & (0.373)        & (0.382)        \\ lexpendB                                &                 & -6.917***      &               & -6.757***      & -6.615***      \\                                         &                 & (0.623)        &               & (0.380)        & (0.379)        \\ prtystrA                                &                 &                & 0.581***      &                & 0.152**        \\                                         &                 &                & (0.121)       &                & (0.062)        \\ Constant                                & 17.806***       & 84.703***      & 21.575***     & 52.039***      & 45.079***      \\                                         & (3.313)         & (3.229)        & (6.121)       & (2.750)        & (3.926)        \\ \hline \multicolumn{1}{c}{Observations}        & 173             & 173            & 173           & 173            & 173            \\ \multicolumn{1}{c}{R2}                  & 0.385           & 0.419          & 0.120         & 0.785          & 0.793          \\ \multicolumn{1}{c}{Adjusted R2}         & 0.382           & 0.416          & 0.114         & 0.783          & 0.789          \\ \multicolumn{1}{c}{Residual Std. Error} & 13.197          & 12.829         & 15.795        & 7.825          & 7.712          \\ \multicolumn{1}{c}{F Statistic}         & 107.238***      & 123.425***     & 23.226***     & 310.694***     & 215.227***     \\ \hline Note:                                   & \multicolumn{5}{l}{*p \textless 0.1     **p \textless 0.05    ***p \textless 0.01    -  Std. Errors in () } \end{tabular} \end{table}

\begin{table}[!htbp] \centering    \caption{\label{tab:three}Matriz de Varianza Covarianza: Modelo 5}    \label{Matriz de Varianza Covarianza}  \begin{tabular}{@{\extracolsep{5pt}} ccccc}  \\[-1.8ex]\hline  \hline \\[-1.8ex]   & Intercept & lexpendA & lexpendB & prtystrA \\  \hline \\[-1.8ex]  Intercept & $15.416$ & $$-$0.395$ & $$-$0.874$ & $$-$0.176$ \\  lexpendA & $$-$0.395$ & $0.146$ & $$-$0.003$ & $$-$0.007$ \\  lexpendB & $$-$0.874$ & $$-$0.003$ & $0.144$ & $0.004$ \\  prtystrA & $$-$0.176$ & $$-$0.007$ & $0.004$ & $0.004$ \\  \hline \\[-1.8ex]  \end{tabular}  \end{table} 

El analista ha estimado un modelo de regresión simple para cada una
de las variables explicativas contra $voteA$ (modelos (1) al (3)),
una especificación en que $voteA$ se estima utilizando como variables
explicativa los gastos de campaña de ambos candidatos (modelo (4))
y finalmente una especificación que incluye, además, la medida de
la popularidad del partido del candidato A (modelo 5) y que corresponde
al modelo general presentado en la ecuación \eqref{eq:Main equation}.
Los resultados de estas estimaciones se presentan en la tabla \ref{tab:two}.
La matriz de varianzas y covarianzas para la estimación del modelo
(5) se presenta en la tabla \ref{tab:three}.\medskip{}

En base a estos resultados:
\begin{enumerate}
\item Comente respecto de los resultados obtenidos por el analista. ¿Tienen
los coeficientes obtenidos los signos esperados?  Explique sus respuestas.

{\color{blue}En efecto, todos los coeficientes obtenidos tienen los signos esperados.
El porcentaje de votos de A debería aumentar con su gasto en campaña y disminuirá con el gasto de sus competidores. Finalmente, también se esperaba que el porcentaje de votos aumentara con la popularidad del partido político del candidato A.   }
\item  ¿Le
parece a usted que la popularidad del partido del candidato A tiene un efecto positivo
sobre porcentaje de votos recibidos por ese candidato?
¿Le parece a usted que esta variable debe ser considerada? ¿Puede
testear formalmente su respuesta? 
Si puede, establezca la forma del
test, incluyendo el estadígrafo de prueba, las hipótesis nula y alternativa,
la distribución muestral del estadígrafo de prueba bajo la nula, sus
grados de libertad y los criterios de no rechazo o rechazo. 

{\color{blue}El coeficiente de $prtystrA$ en el modelo 5 es positivo y significativo
al 5\%. El test t que aparece en la tabla (2) tiene como hipótesis alternativa que $\beta_{prtystrA}$ es distinta de 0. Debemos chequear que es positivo, por lo tanto, deberíamos realizar un test de una cola:  


$H_{0}:\beta_{prtystrA}=0$ \\ $H_{1}: \beta_{prtystrA}>0$

$\hat{t}=\frac{\hat{\beta}_{prtystrA}-0}{\sqrt{\hat{\sigma}_{prtystrA}^{2}}}\underset{H_{0}}{\sim}t_{173-4}\Rightarrow\hat{t}=\frac{0.152}{0.062}=2.45$ 

% \begin{figure}[h!]
%     \centering
%     \includegraphics[width=0.5\linewidth]{guias/test_t.png}
% \end{figure}

Por ende, $\hat{\beta}_{prtystrA}$ es significativo al 1\% en un
test unilateral y signficativo al 5\% en un test bilateral. El test debería ser unilateral, ya que se desea testear si el efecto de la popularidad de un partido es positivo. }


\item ¿Cuál es la interpretación del coeficiente $\beta_{1}$ (intercepto
o constante)? Al analista le preocupa el hecho de que el estimador
de la constante es muy diferente entre el modelo (1) y los modelos
(4) y (5). ¿Puede usted explicar por qué sucede esto?

{\color{blue}

La constante captura todo el efecto promedio de \(Y\) que no es explicado por los regresores.

Si las variables no están centradas (es decir, sus medias no son 0), cualquier cambio en el conjunto de variables independientes puede alterar cómo se calcula el promedio de \(Y\) para \(X=0\), afectando directamente el valor de la constante estimada.



\[
\hat{\beta_0} = \bar{Y} - \hat{\beta_1} \bar{X_1}
\]

Al incluir una nueva variable \(X_2\), la constante estimada ahora será:

\[
\hat{\beta_0} = \bar{Y} - \hat{\beta_1 }\bar{X_1} - \hat{\beta_2} \bar{X_2}
\]

Por lo tanto, esperamos ver cambios en la constante al incluir o sacar variables explictivas. }




\item Utilizando los resultados del modelo (5), testee la siguiente afirmación:
Un aumento del 1\% en los gastos del candidato A es compensado por
un incremento del 1\% en los gastos del candidato B. Establezca la
forma del test, incluyendo el estadígrafo de prueba, las hipótesis
nula y alternativa, la distribución muestral del estadígrafo de prueba
bajo la nula, sus grados de libertad y los criterios de no rechazo
o rechazo. 

{\color{blue}$H_{0}:\beta_{1}+\beta_{2}=0$ , $H_{1}:\beta_{1}+\beta_{2}\neq0$

$\hat{t}=\frac{\hat{\beta}_{1}+\hat{\beta}_{2}-0}{\sqrt{\hat{var}\left(\hat{\beta}_{1}+\hat{\beta}_{2}\right)}}=\frac{\hat{\beta}_{1}+\hat{\beta}_{2}}{\sqrt{\hat{var}\left(\hat{\beta}_{1}\right)+\hat{var}\left(\hat{\beta}_{2}\right)+2\hat{cov}\left(\hat{\beta}_{1},\hat{\beta}_{2}\right)}}\underset{H_{0}}{\sim}t_{173-4}$

$\Rightarrow\hat{t}=\frac{6.083-6.615}{\sqrt{0.146+0.144-2\cdot0.003}}=-0.9983$


\begin{figure}[h!]
    \centering
    \includegraphics[width=0.9\linewidth]{Ejercicios//ejercicios_inferencia/t169.png}
\end{figure}
A los niveles de significancia del 10\%, 5\% y 1\%, el estadístico de prueba no se encuentra dentro de la región de rechazo. Por lo tanto, no se rechaza la hipótesis nula incluso al nivel del 10\%.}


\item En base a los resultados obtenidos para el modelo (5), lleve a cabo
un test de significancia global del modelo. En su respuesta, establezca la forma del test, incluyendo
el estadígrafo de prueba, las hipótesis nula y alternativa, la distribución
muestral del estadígrafo de prueba bajo la nula, sus grados de libertad
y los criterios de no rechazo o rechazo.

{\color{blue}$H_{0}:\beta_{2}=\beta_{3}=\beta_{4}=0$ \\
$H_{1}:$ Al menos uno
es distinto de cero.

 $\hat{F}=\frac{\frac{R^{2}}{k}}{\frac{1-R^{2}}{n-k-1}}\underset{H_{0}}{\sim}F_{k},_{n-k-1}\Rightarrow\hat{F}=\frac{\frac{0.793}{3}}{\frac{1-0.793}{169}}=215.8$

\begin{figure}
    \centering
    \includegraphics[width=0.9\linewidth]{Ejercicios//ejercicios_inferencia/F_3_169.png}
\end{figure}
 Por ende, se rechaza la nula a cualquier nivel estándar de significancia.}

\end{enumerate}

\end{document}